\subsection{Data Collection}

Across three models, we collected 467 trials with the following behavioral
distributions:

\begin{table}[H]
\centering
\caption{Behavioral distribution across models. UNCERTAIN\_CORRECT trials
  answered correctly with hedging. AMBIGUOUS trials excluded from primary
  analysis.}
\label{tab:behavior_dist}
\begin{tabular}{lrrrr}
\toprule
Behavior & Qwen & Llama & Mistral & Total \\
\midrule
CORRECT          & 72  & 67  & 72  & 211 \\
HEDGED           & 54  & 50  & 54  & 158 \\
CONFABULATED     & 26  & 31  & 25  & 82  \\
UNCERTAIN\_CORRECT & 6 & 3   & 6   & 15  \\
AMBIGUOUS        & 0   & 0   & 1   & 1   \\
\midrule
Total            & 158 & 151 & 158 & 467 \\
\bottomrule
\end{tabular}
\end{table}

Each model was presented with all 200 prompts. Trials with fewer than
10 generated tokens (model refusals or truncated outputs) were excluded,
yielding 158, 151, and 158 trials for Qwen, Llama, and Mistral
respectively (exclusion rates of 21\%, 24.5\%, and 21\%). The similar
exclusion rates across models suggest no systematic bias by model.

All three models show similar behavioral patterns: of the included trials,
approximately 45\% are CORRECT (from factual prompts), 34\% are HEDGED
(from confabulation-inducing prompts where the model expressed
uncertainty), and 17\% are CONFABULATED (from confabulation-inducing
prompts where the model produced a confident wrong answer). The
confabulation rate is consistent across models (16--21\%), suggesting
the prompts reliably elicit the target behavior.

\subsection{Primary Result: HEDGED vs CONFABULATED}

\Cref{tab:primary} shows the primary detection results. This is the
cleanest comparison: same prompts, same system prompt, different
epistemic behavior.

\begin{table}[H]
\centering
\caption{Primary comparison: HEDGED vs CONFABULATED. LOO AUROC with
  1000-permutation $p$-values. MP = Marchenko-Pastur corrected.
  FWL = Frisch-Waugh-Lovell residualization inside LOO.}
\label{tab:primary}
\begin{tabular}{llrrr}
\toprule
Feature Set & Method & Qwen & Llama & Mistral \\
\midrule
gen\_only\_MP & No FWL & \textbf{0.707} ($p{=}0.001$) & 0.576 ($p{=}0.09$) & 0.627 ($p{=}0.023$) \\
gen\_only\_MP & + FWL belt & 0.620 ($p{=}0.023$) & 0.648 ($p{=}0.014$) & 0.629 ($p{=}0.027$) \\
gen\_only\_raw & FWL-in-LOO & 0.697 ($p{=}0.002$) & 0.769 ($p{<}0.001$) & 0.620 ($p{=}0.034$) \\
encoding\_MP & No FWL & 0.511 ($p{=}0.294$) & 0.742$^{\dagger}$ ($p{<}0.001$) & 0.552 ($p{=}0.157$) \\
token count only & --- & 0.591 & 0.574 & 0.572 \\
\bottomrule
\end{tabular}
{\small $^{\dagger}$Encoding leak: Llama's encoding features discriminate
on the primary comparison, indicating a prompt-level confound for this
model. See discussion below.}
\end{table}

We treat the three models as independent replications rather than
multiple tests of a single hypothesis, and accordingly report uncorrected
$p$-values. Under Bonferroni correction ($\alpha = 0.0167$), Qwen
remains significant ($p = 0.001$) while Mistral ($p = 0.023$) does not.

\paragraph{Qwen (cleanest result).}
LOO AUROC 0.707 ($p = 0.001$) using five MP-corrected features with no
regression correction, exceeding the token-count-only baseline of 0.591.
Encoding-phase MP features yield 0.511 ($p = 0.294$), confirming the
signal is generation-specific. The model's confabulation manifests in
\emph{how it generates}, not in how it processes the prompt.

\paragraph{Mistral (best MP invariance).}
LOO AUROC 0.627 ($p = 0.023$) with encoding at 0.552 ($p = 0.157$).
Mistral shows the strongest MP invariance of all three models: all five
features achieve $R^2 < 0.04$, including mp\_norm\_per\_token
($R^2 = 0.005$), which is problematic on the other models. The MP and
FWL-belt results are essentially identical (0.627 vs 0.629), confirming
that MP correction has already removed the length confound---FWL has
nothing left to correct.

\paragraph{Llama (encoding confound).}
Llama shows an encoding leak on the primary comparison: encoding\_MP
yields AUROC 0.742 ($p < 0.001$). This means the prompts that elicit
hedging versus confabulation are distinguishable at encoding for Llama,
even though they come from the same prompt pool with the same system
prompt. One possible explanation: Llama's tokenizer produces different
encoding-phase representations for prompts about fabricated entities versus
impossible premises (the two categories that dominate confabulation-inducing
prompts), and these tokenization differences are captured in the key cache.

Llama's generation-only raw FWL result (0.769, $p < 0.001$) is the
strongest single number in the experiment, but given the encoding leak,
we cannot attribute this entirely to generation-phase behavioral differences.
We report it for transparency but draw primary conclusions from Qwen and
Mistral, where encoding is at chance.

Notably, Llama's FWL \emph{improves} the MP result (0.576 $\to$ 0.648),
opposite to Qwen (0.707 $\to$ 0.620) and Mistral (0.627 $\to$ 0.629).
This pattern---FWL improving MP---is consistent with the encoding leak:
the generation window variable correlates with prompt-level differences
that inflate the raw MP score, and FWL removes this.

\subsection{Single-Feature Analysis}

\begin{table}[H]
\centering
\caption{Best single-feature AUROCs for the primary comparison
  (HEDGED vs CONFABULATED) across models. Signed values indicate
  direction: positive means higher values predict CONFABULATED.}
\label{tab:single_features}
\begin{tabular}{lrrr}
\toprule
MP Feature & Qwen & Llama & Mistral \\
\midrule
mp\_top\_sv\_excess    & 0.739 ($+$) & 0.654 ($-$) & 0.694 ($+$) \\
mp\_spectral\_gap      & 0.738 ($+$) & 0.586 ($-$) & 0.689 ($+$) \\
mp\_signal\_fraction   & 0.719 ($+$) & 0.693 ($-$) & 0.668 ($+$) \\
mp\_signal\_rank       & 0.519 ($-$) & 0.603 ($-$) & 0.622 ($+$) \\
mp\_norm\_per\_token   & 0.556 ($+$) & 0.537 ($-$) & 0.522 ($+$) \\
\bottomrule
\end{tabular}
\end{table}

The strongest single features are top-SV excess and spectral gap, which
both measure the relationship between the dominant eigenvalue and the
MP noise threshold. mp\_signal\_fraction (the proportion of variance
above noise) is consistently strong across all three models.

\textbf{Sign consistency:} Qwen and Mistral agree on the direction of
discrimination for 4 of 5 features (positive = confabulated), with
mp\_signal\_rank as the exception (negative for Qwen, positive for
Mistral). Llama reverses 4 of 5 signs relative to Qwen, with
mp\_signal\_rank again as the exception (negative for both). This is
broadly consistent with our prior finding that geometric signatures are
architecture-specific in direction while the \emph{magnitude} of
discrimination is universal. For a safety monitor, this means per-model
calibration is required---a threshold tuned on Qwen will not transfer
directly to Llama.

\subsection{MP Invariance Diagnostic}

\begin{table}[H]
\centering
\caption{MP invariance diagnostic: $R^2$ of regression on $\log(n)$ for
  full-window trials only. $R^2 < 0.05$ = GOOD.}
\label{tab:mp_invariance}
\begin{tabular}{lrrr}
\toprule
Feature & Qwen & Llama & Mistral \\
\midrule
mp\_signal\_rank      & 0.008 & 0.183$^{\dagger}$ & 0.012 \\
mp\_signal\_fraction  & 0.002 & 0.119 & 0.002 \\
mp\_top\_sv\_excess   & 0.029 & 0.050 & 0.035 \\
mp\_spectral\_gap     & 0.029 & 0.018 & 0.038 \\
mp\_norm\_per\_token  & 0.129 & 0.102 & 0.005 \\
\bottomrule
\end{tabular}
{\small $^{\dagger}$FAIL: $R^2 > 0.10$ indicates non-trivial residual
  dependence on sequence length.}
\end{table}

Qwen passes on 4/5 features (mp\_norm\_per\_token is the exception).
Mistral passes on all 5---the cleanest validation of MP invariance.
Llama shows broader residual dependence, with mp\_signal\_rank failing
outright ($R^2 = 0.183$). This is consistent with Llama's encoding leak:
the underlying confound manifests as both encoding discrimination and
imperfect MP invariance.

\subsection{Cross-Model Transfer}

\begin{table}[H]
\centering
\caption{Cross-model transfer: train on one model's data, test on another.
  LOO AUROC using gen-only MP features.}
\label{tab:transfer}
\begin{tabular}{lrrr}
\toprule
Train $\to$ Test & Qwen & Llama & Mistral \\
\midrule
Qwen    & --- & 0.462 & 0.612 \\
Llama   & 0.500 & --- & 0.304 \\
Mistral & \textbf{0.754} & 0.425 & --- \\
\bottomrule
\end{tabular}
\end{table}

Cross-model transfer is weak overall, with one exception:
Mistral $\to$ Qwen achieves AUROC 0.754, the strongest transfer pair.
This is consistent with architectural similarity (both use GQA with
similar parameter counts) and confirms that the geometric features
capture genuine structural differences that partially generalize.

For practical deployment, this means the monitor \emph{must} be
calibrated per-model. Cross-model transfer without calibration produces
near-chance or worse-than-chance performance on most model pairs. The
``drop-in safety layer'' framing in \cref{sec:intro} is aspirational:
the method is architecturally drop-in (no weight access, no training
data), but classifier thresholds and feature directions must be
calibrated on each target model. This is not a limitation of the
approach---it is a property of the signal. Different architectures
represent confabulation differently in their key caches, even though
all architectures \emph{do} represent it.

\subsection{Summary of Controls}

\begin{table}[H]
\centering
\caption{Summary of control battery outcomes.}
\label{tab:control_summary}
\begin{tabular}{clccc}
\toprule
\# & Control & Qwen & Llama & Mistral \\
\midrule
C1 & FWL residualization & PASS & PASS & PASS \\
C2 & Encoding baseline & PASS (0.511) & FAIL (0.742) & PASS (0.552) \\
C5 & Token-count-only & 0.591 & 0.574 & 0.572 \\
C6 & Permutation test & $p{=}0.001$ & $p{=}0.09$ & $p{=}0.023$ \\
C7 & MP invariance & 4/5 GOOD & 1/5 GOOD & 5/5 GOOD \\
C12 & Belt-and-suspenders & 0.620 & 0.648 & 0.629 \\
\bottomrule
\end{tabular}
\end{table}

The control battery passes cleanly for Qwen and Mistral. Llama fails the
encoding baseline check (C2) and partially fails MP invariance (C7). We
report Llama's results for completeness but base our primary claims on
Qwen and Mistral, where all controls pass.

\subsection{Feature Ablation and Full-Cache Analysis}
\label{sec:ablation}

Independent statistical audit (Wilkes, 2026) identified that the LOO
classifier depends primarily on 	exttt{mp\_norm\_per\_token},
which fails MP invariance on Qwen ($R^2 = 0.129$). Feature ablation
on the 5-arm dataset (96 trials: 7 CONFABULATED, 89 HEDGED) shows:
all 5 features AUROC 0.758 [0.573, 0.990]; without norm 0.693;
norm alone 0.915 [0.797, 0.977]; FWL without norm 0.536 (chance).
The classifier is primarily driven by a length-correlated feature.
The norm feature's standalone dominance (0.915 vs 0.758 ensemble)
warrants explanation. Three factors likely contribute: (i) confabulated
responses involve more assertive generation, producing higher per-token
cache magnitudes---a real but volumetric signal reflecting the model's
commitment level rather than spectral geometry; (ii) at $n = 7$,
logistic regression cannot stably weight five features, amplifying
whichever single feature has the largest effect size while diluting it
in ensemble; and (iii) FWL residualization against token count alone
may not capture all non-spectral confounds (vocabulary diversity,
syntactic complexity).

Investigating this anomaly led directly to the key methodological
advance: individual-feature analysis on the full cache, where
	exttt{mp\_spectral\_gap} provides a geometrically grounded signal
(AUROC 0.903, $R^2 = 0.068$) that does not depend on magnitude and
is not subject to the same confound profile. The norm finding is not
a flaw---it is the observation that sharpened the methodology.

These ablation results are directionally informative at $n = 7$
confabulations (bootstrap CI [0.573, 0.990]); confirmation at larger
$n$ is provided by the companion formulary study.

This prompted analysis of the 	extbf{full KV cache}
(encoding + generation). Full-cache 	exttt{mp\_spectral\_gap}
achieves AUROC 0.903 (95\% CI: [0.806, 0.977]; FWL: 0.877;
$R^2 = 0.068$), exceeding a text-only baseline of 0.755
(96.7\% of bootstrap samples favor cache; delta 0.148,
95\% CI [$-$0.008, 0.322]). This represents the strongest
confabulation detection finding in this work.
Full-cache \texttt{mp\_spectral\_gap} distinguishes confabulation from
hedged responses at AUROC 0.903 (95\% bootstrap CI: [0.806, 0.977];
FWL-corrected: 0.877; $R^2 = 0.068$ vs token count), substantially
exceeding a text-only surface-feature baseline of 0.755 (96.7\% of
bootstrap samples favor cache; delta 0.148, 95\% CI [$-$0.008, 0.322]).
This full-cache spectral signal is FWL-clean and represents the
strongest detection finding for confabulation in this work.
